% Architecture SPDnet, dessinee de haut en bas: c'est le seul sens qui laisse
% aux boites une taille lisible dans les 355pt du texte.
% Utilise dans la \cref{sec:spdnet-couches}.
% Styles et briques communs aux figures du \cref{ch:spdnet}.
% Les teintes reprennent celles de gfx/tikZ/riemann/style.tex pour que les
% figures des deux chapitres se lisent comme une meme famille.
% Le texte du memoire ne fait que 355pt (12,5 cm): toutes les figures de ce
% chapitre s'ecoulent donc de haut en bas, jamais de gauche a droite.
% Ce fichier peut etre \input plusieurs fois.
\definecolor{slayer}{RGB}{90,90,90}      % boite de couche, gris neutre
\definecolor{sspd}{RGB}{150,40,40}       % ce qui vit sur le cone SPD
\definecolor{sstiefel}{RGB}{40,70,140}   % ce qui vit sur Stiefel
\definecolor{seucl}{RGB}{50,110,120}     % ce qui vit dans un espace euclidien
\definecolor{sback}{RGB}{175,95,45}      % retropropagation
\tikzset{
  spdnet/.style={
    >={Stealth[length=5pt,width=4pt]},
    line join=round,
    font=\small,
    node distance=5mm,
  },
  % Boite de couche generique. La largeur est fixee pour que la colonne soit
  % alignee et que le texte respire.
  slayerbox/.style={
    draw=slayer!70, thick, rounded corners=2pt, align=center,
    inner sep=4pt, minimum height=9mm, minimum width=42mm, fill=white,
  },
  sstiefelbox/.style={slayerbox, draw=sstiefel!75, fill=sstiefel!7},
  splainbox/.style={slayerbox, fill=slayer!5},
  seuclbox/.style={slayerbox, draw=seucl!75, fill=seucl!7},
  smeanbox/.style={slayerbox, draw=sspd!75, fill=sspd!9},
  % Boite etroite, pour les colonnes doubles.
  snarrow/.style={slayerbox, minimum width=32mm},
  sarrow/.style={->, thick, slayer!85},
  sbackarrow/.style={->, thick, sback, dashed},
  sblock/.style={draw=slayer!45, thick, dashed, rounded corners=3pt, inner sep=5pt},
  slabel/.style={font=\footnotesize, inner sep=2pt},
  stiny/.style={font=\scriptsize, inner sep=1.5pt},
  % Accolade de regroupement a droite ou a gauche d'une colonne.
  sbrace/.style={decorate, decoration={brace, amplitude=4pt}, slayer!70},
}

\begin{tikzpicture}[spdnet]

  \node[splainbox] (in) {entrée $\MCov_0\in\SPD{\dlayer{0}}$};

  \node[sstiefelbox, below=of in] (bimap1) {\BiMap{} \ $\Wlayer_1\in\Stiefel{\dlayer{0}}{\dlayer{1}}$};
  \node[splainbox, below=of bimap1] (reeig1) {\ReEig{} \ seuil $\reeigeps$};
  \node[smeanbox, below=of reeig1] (bn1) {\BatchNorm{} \ $\Fmean,\ \Covbias_1$};

  \node[below=6mm of bn1] (dots) {$\vdots$};

  \node[sstiefelbox, below=6mm of dots] (bimapL) {\BiMap{} \ $\Wlayer_\Nlayer\in\Stiefel{\dlayer{\Nlayer-1}}{\dlayer{\Nlayer}}$};
  \node[splainbox, below=of bimapL] (reeigL) {\ReEig{} \ seuil $\reeigeps$};

  \node[splainbox, below=8mm of reeigL] (logeig) {\LogEig};
  \node[seuclbox, below=of logeig] (fc) {$\softmax$ \ $\MSoftmax,\Vsoftmax$};
  \node[splainbox, below=of fc] (out) {sortie $\hat{\classlabel}\in\R^{\Nclass}$};

  \foreach \a/\b in {in/bimap1, bimap1/reeig1, reeig1/bn1, bn1/dots, dots/bimapL,
                     bimapL/reeigL, reeigL/logeig, logeig/fc, fc/out}
    \draw[sarrow] (\a) -- (\b);
  \coordinate (mid) at ($(logeig.south)!0.5!(fc.north)$);
  \node[stiny, slayer, anchor=west] at ([xshift=2mm]mid) {$\vech$};

  % Regroupement des deux blocs BiRe.
  \node[sblock, fit=(bimap1)(bn1)] (b1) {};
  \node[stiny, slayer, rotate=90, anchor=south] at ($(b1.west)+(-1mm,0)$) {bloc \BiRe{} 1};
  \node[sblock, fit=(bimapL)(reeigL)] (bL) {};
  \node[stiny, slayer, rotate=90, anchor=south] at ($(bL.west)+(-1mm,0)$) {bloc \BiRe{} $\Nlayer$};

  % Colonne de droite: ou vivent les objets.
  \coordinate (x) at ([xshift=6mm]in.east);
  \draw[sspd!75, thick] ([yshift=1mm]in.north-|x) -- ([yshift=-1mm]logeig.north-|x);
  \node[stiny, sspd, anchor=west, align=left, text width=32mm]
    at ($([yshift=1mm]in.north-|x)!0.5!([yshift=-1mm]logeig.north-|x)+(2mm,0)$)
    {les représentations restent dans le cône $\SPD{\dlayer{\ilayer}}$};
  \draw[seucl!75, thick] ([yshift=-1mm]logeig.north-|x) -- ([yshift=-1mm]out.south-|x);
  \node[stiny, seucl, anchor=west, align=left, text width=32mm]
    at ($([yshift=-1mm]logeig.north-|x)!0.5!([yshift=-1mm]out.south-|x)+(2mm,0)$)
    {espace tangent en $\Id{\dlayer{\Nlayer}}$, puis $\R^{\Nclass}$};

\end{tikzpicture}
